% !TEX root = ../SANER2027-main-labeling.tex
% Discussion: failure analysis and implications.

\section{Discussion}
\label{sec:discussion}

\subsection{Failure Analysis}
\label{sec:disc-failures}

We manually analyzed the errors of Qwen-32B at each method's best configuration. We randomly sampled 50 misclassified issues each from \rag, \frag, and LoRA fine-tuning, 150 in total. Each method's sample has 30 true-question, 10 true-bug, and 10 true-feature errors, which mirrors the label distribution of the errors pooled over the three methods. Two authors analyzed each case independently and resolved disagreements through discussion. We found three categories.

\myparagraph{Wrong template (116 of 150, 77\%)} The issue's content and label contradict the template it was filed under. All of these errors are questions and feature requests filed under a bug-report template (\textit{Steps to reproduce}, \textit{Expected}, \textit{Actual}, \textit{System Information}), whose structure leads the LLM to predict bug regardless of the content. For example, \texttt{vscode\#193985}\footnote{\url{https://github.com/microsoft/vscode/issues/193985}} is a support question that opens with a literal \texttt{Type:Bug} because it uses VS Code's bug template.

\myparagraph{Hybrid intent (20 of 150, 13\%)} The issue combines intents and could defensibly carry more than one label, mostly because a fix or a feature request is layered into a question or bug report. For example, in \texttt{kubernetes\#120364}\footnote{\url{https://github.com/kubernetes/kubernetes/issues/120364}} the author reports a concrete bug and, in the same body, requests a new feature: ``\textit{Please, give us multiarchitecture image! and fix monoarchitecture images too!}''

\myparagraph{Label noise (14 of 150, 9\%)} The ground-truth label does not match the content. For example, \texttt{tensorflow\#61710}\footnote{\url{https://github.com/tensorflow/tensorflow/issues/61710}}, titled ``\textit{Will there be an MLP model in the future version?}'', requests a new API feature but is labeled as a bug.

Wrong templates affect all three approaches. We hypothesize that the headings of a bug template are a strong cue for the bug label, both for the retriever, which then returns bug-labeled neighbors, and for the LLM itself. The other two categories are harder for any text-based classifier. For the 34 errors with hybrid intent or label noise (23\%), another label is defensible or the ground truth is wrong, so they limit the macro $F_1$ that any classifier can reach on this benchmark.

We also sampled 30 invalid outputs each from \rag\ and \frag. In 44 of the 60 cases (73\%), the model continued the issue body or its own reasoning. In 14 (23\%), it produced degenerate text such as repeated digits and punctuation. In the remaining 2 (3\%), it gave a label outside the three, such as \texttt{documentation}.

\subsection{Implications}
\label{sec:disc-implications}

\myparagraph{Labeled data and retraining} With \frag, a project's own labeled issues give, without training and with 33\% less peak GPU memory, a macro $F_1$ within 1.0 point of fine-tuning on $11\times$ as many issues pooled from eleven projects; the difference is not statistically significant, while \rag\ trails by 2.8 points (\Cref{sec:rq3}). The comparison is starker for a project that cannot draw on other projects' labeled issues: fine-tuned on the same issues, the model scores 3.8--6.0 points below \frag. Building the index for a project needs only its labeled issues and one embedding pass, whereas pooled fine-tuning needs labeled issues collected from other projects. A new labeled issue then becomes an index update rather than a retraining run, so the examples can follow a project as its issues change, and the LLM can be replaced by a newer one without retraining.

\myparagraph{Where fine-tuning keeps an edge} Fine-tuning keeps three advantages. It is significantly ahead at Qwen-7B and Qwen-14B, the sizes at which pooling the eleven projects' issues helps it most (\Cref{sec:rq3}). Its inference is faster, most clearly at Qwen-32B, where \frag\ takes longer in total even when fine-tuning's training is included; this advantage grows with the number of issues labeled between two training runs. And it almost never produces invalid outputs, whereas the RAG variants need the \knn\ fallback to give every issue a valid label.

\myparagraph{Error profiles} The approaches differ less in macro $F_1$ than in which errors they make. Fine-tuning finds more bugs at every size, whereas \frag\ raises fewer false bug alarms at three of the four sizes. Macro $F_1$ weighs both errors equally, but their costs need not be equal: a false bug label costs a triager's time, whereas a missed bug can delay a fix. Neither approach removes the question-to-bug confusion.

\myparagraph{Question-to-bug confusion} The confusion persists under every method we test, including fine-tuning, as in prior \irc\ work (\Cref{sec:02_related}). The filter reduces it without training but cannot remove it. Many of the remaining errors are questions filed under a bug template (\Cref{sec:disc-failures}), which read like bug reports to the retriever and the LLM alike.

\myparagraph{Retrieval alone as a floor} Without an LLM, \knn\ reaches 59.5\% macro $F_1$, 1.8 points below Qwen-3B's zero-shot prompting, with 0.2 instead of 2.9~GB of GPU memory. Neither the neighbors' labels nor the LLM alone reaches the scores of \rag, which outperforms both at every size (\Cref{sec:rq1}); the gain comes from giving the LLM the labeled neighbors.

\myparagraph{Classification and templates} Our failure analysis shows that users often file support questions and feature requests under the bug-report template, which leads the classifier to predict bug regardless of the content. Most of the sampled errors are of this kind, so correcting the template when an issue is filed could prevent many of them. This matters for the LLM-assisted triage now offered on platforms such as GitHub~\cite{githubcopilot} and GitLab~\cite{gitlabduo}, where the reporter describes the issue to an LLM that suggests a template and labels. Report validation matters as well, since mislabeled and duplicate reports are common~\cite{antoniol2008bug, pandey2017automated}. Prior studies validate bug reports~\cite{khatib2025assertflip, wang2024aegis, ahmed2025otter, dincc2025judge}, but combining classification and validation in one workflow remains underexplored.
